% \appendix is invoked in main.tex before all appendix inputs.

\section{Participant Recruitment and Sample Construction}
\label{app:participants}

\subsection{Sampling rationale}

The study is designed to examine workflow ownership without assuming skill
engineering expertise. We therefore recruit for high \emph{workflow
accountability} and variation in \emph{artifact familiarity}. Workflow
accountability means that the participant defines task goals, judges the
acceptability of agent behavior, depends on the workflow's output, or bears
responsibility for its consequences. Artifact familiarity concerns the extent
to which the participant can inspect, author, or modify reusable instructions.
We target 12--16 completed sessions, with a planned stopping point of 14 unless
attrition or an imbalanced sample requires up to two additional participants.
This target prioritizes information-rich variation. Statistical
representativeness lies outside the study's purpose.

We seek variation along the following dimensions:

\begin{table}[h]
\centering
\caption{Maximum-variation sampling dimensions. These are sampling targets,
not experimental conditions.}
\label{tab:sampling-dimensions}
\begin{tabularx}{\linewidth}{P{0.25\linewidth}YY}
\toprule
\textbf{Dimension} & \textbf{Lower end} & \textbf{Higher end} \\
\midrule
Artifact familiarity & Rarely inspects or edits reusable instructions; may ask a teammate or AI to make changes & Regularly authors, inherits, reviews, or versions system prompts, routines, or skills \\
Workflow consequence & Primarily time, quality, or inconvenience costs & Material financial, coordination, reputational, or operational consequences, excluding regulated high-stakes decisions \\
Workflow structure & Primarily generates or transforms information & Uses multiple tools, priorities, confirmations, or conditional actions \\
Ownership arrangement & Individual workflow & Shared, inherited, or organizational workflow \\
Domain & Research, writing, analysis, or personal productivity & Product, operations, software, customer support, or other organizational work \\
\bottomrule
\end{tabularx}
\end{table}

\subsection{Inclusion and exclusion criteria}

Participants must satisfy all inclusion criteria:

\begin{enumerate}[leftmargin=*]
    \item They have used an LLM assistant or agent for the same recurring family
    of tasks at least five times in the previous three months.
    \item The workflow is shaped by a persistent instruction or configuration
    artifact, such as custom instructions, a saved system prompt, a reusable
    template, a routine, or an agent skill. The participant need not have
    authored it, but must be able to describe its role and provide at least a
    redacted excerpt, reconstruction, or screen-shared view.
    \item They can identify an execution within the previous six months that
    they considered unacceptable, surprising, or unsafe to reuse without
    review.
    \item They had responsibility for judging the workflow's goals or outcomes,
    even if another person or AI implemented the configuration.
    \item They can provide a redacted failure artifact and either a typical
    successful execution or a description of expected behavior.
\end{enumerate}

We exclude workflows whose discussion would require disclosure of protected
health information, active legal representation, classified material,
credentials, or other information that cannot be safely redacted. We also
exclude one-off prompting, workflows without any persistent configuration, and
participants who cannot describe a concrete breakdown. The study examines
repair reasoning, not the correctness of decisions in regulated domains.

\subsection{Recruitment screener}

The screener should be administered before scheduling. Bracketed response
options are suggestions and may be adapted to the recruitment platform.

\begin{enumerate}[leftmargin=*]
    \item Describe one task for which you repeatedly use the same AI assistant,
    agent, custom GPT, saved prompt, template, routine, or skill.
    \item Approximately how many times have you used this workflow in the last
    three months? [1--2; 3--4; 5--10; 11--25; more than 25]
    \item Who decides whether the workflow's output or action is acceptable?
    [Primarily me; shared; primarily another person]
    \item Who can inspect or change the reusable instructions? [Me; a teammate;
    AI on my behalf; platform owner; unknown]
    \item Have you encountered a concrete execution that caused you to modify,
    consider modifying, reduce reliance on, or stop using this workflow?
    [Yes/No; brief description]
    \item What could go wrong if the workflow behaves incorrectly? [Time loss;
    poor quality; financial cost; coordination cost; reputational cost; other]
    \item Can you provide a redacted screenshot, output, trace, instruction
    excerpt, version diff, or reconstruction related to that incident? [Yes/No]
    \item Does the material contain information that cannot be safely redacted
    or discussed in a research session? [Yes/No]
\end{enumerate}

\subsection{Participant table template}

Populate Table~\ref{tab:participant-template} only after the sample is final.
Use broad role and domain labels to avoid re-identification.

\begin{table*}[t]
\centering
\caption{Participant characteristics. Replace all placeholders with verified,
de-identified information.}
\label{tab:participant-template}
\begin{tabularx}{\textwidth}{P{0.045\textwidth}P{0.13\textwidth}P{0.22\textwidth}P{0.12\textwidth}P{0.14\textwidth}Y}
\toprule
\textbf{ID} & \textbf{Role / domain} & \textbf{Recurring workflow} & \textbf{Artifact access} & \textbf{Artifact familiarity} & \textbf{Primary consequence} \\
\midrule
P1 & \DataNeeded{ROLE} & \DataNeeded{WORKFLOW} & \DataNeeded{ACCESS} & \DataNeeded{LEVEL} & \DataNeeded{CONSEQUENCE} \\
P2 & \DataNeeded{ROLE} & \DataNeeded{WORKFLOW} & \DataNeeded{ACCESS} & \DataNeeded{LEVEL} & \DataNeeded{CONSEQUENCE} \\
P3 & \DataNeeded{ROLE} & \DataNeeded{WORKFLOW} & \DataNeeded{ACCESS} & \DataNeeded{LEVEL} & \DataNeeded{CONSEQUENCE} \\
P4 & \DataNeeded{ROLE} & \DataNeeded{WORKFLOW} & \DataNeeded{ACCESS} & \DataNeeded{LEVEL} & \DataNeeded{CONSEQUENCE} \\
P5 & \DataNeeded{ROLE} & \DataNeeded{WORKFLOW} & \DataNeeded{ACCESS} & \DataNeeded{LEVEL} & \DataNeeded{CONSEQUENCE} \\
P6 & \DataNeeded{ROLE} & \DataNeeded{WORKFLOW} & \DataNeeded{ACCESS} & \DataNeeded{LEVEL} & \DataNeeded{CONSEQUENCE} \\
P7 & \DataNeeded{ROLE} & \DataNeeded{WORKFLOW} & \DataNeeded{ACCESS} & \DataNeeded{LEVEL} & \DataNeeded{CONSEQUENCE} \\
P8 & \DataNeeded{ROLE} & \DataNeeded{WORKFLOW} & \DataNeeded{ACCESS} & \DataNeeded{LEVEL} & \DataNeeded{CONSEQUENCE} \\
P9 & \DataNeeded{ROLE} & \DataNeeded{WORKFLOW} & \DataNeeded{ACCESS} & \DataNeeded{LEVEL} & \DataNeeded{CONSEQUENCE} \\
P10 & \DataNeeded{ROLE} & \DataNeeded{WORKFLOW} & \DataNeeded{ACCESS} & \DataNeeded{LEVEL} & \DataNeeded{CONSEQUENCE} \\
P11 & \DataNeeded{ROLE} & \DataNeeded{WORKFLOW} & \DataNeeded{ACCESS} & \DataNeeded{LEVEL} & \DataNeeded{CONSEQUENCE} \\
P12 & \DataNeeded{ROLE} & \DataNeeded{WORKFLOW} & \DataNeeded{ACCESS} & \DataNeeded{LEVEL} & \DataNeeded{CONSEQUENCE} \\
P13 & \DataNeeded{ROLE} & \DataNeeded{WORKFLOW} & \DataNeeded{ACCESS} & \DataNeeded{LEVEL} & \DataNeeded{CONSEQUENCE} \\
P14 & \DataNeeded{ROLE} & \DataNeeded{WORKFLOW} & \DataNeeded{ACCESS} & \DataNeeded{LEVEL} & \DataNeeded{CONSEQUENCE} \\
\bottomrule
\end{tabularx}
\end{table*}

\section{Pre-Study Artifact Collection and Data Handling}
\label{app:artifacts}

\subsection{Artifact request}

Two to three days before the session, participants receive the following
request:

\begin{quote}
Please choose one recurring task for which you use an AI assistant or agent with
a persistent configuration, such as custom instructions, a saved prompt,
template, routine, or skill. Provide, if available: (1) one execution that you
considered unacceptable or surprising; (2) one typical or successful execution
from the same workflow; (3) the relevant reusable instructions or a redacted
excerpt; and (4) any note, version history, or change you considered after the
incident. You may redact names, organizations, values, and proprietary details,
reconstruct the case with equivalent placeholders, or show the materials only
through screen sharing during the session.
\end{quote}

Participants also answer a brief artifact inventory:

\begin{enumerate}[leftmargin=*]
    \item What recurring task was the workflow intended to support?
    \item Who created and who currently maintains the reusable configuration?
    \item Which parts of the configuration can you inspect or edit?
    \item What happened in the selected breakdown?
    \item What did you do afterward: rerun, edit, ask someone else to edit,
    manually take over, rollback, reduce automation, or stop using the workflow?
    \item What behavior did you hope would remain unchanged?
    \item What evidence, if any, convinced you that the incident was resolved?
\end{enumerate}

\subsection{Privacy and data minimization}

Participants are instructed not to submit credentials, private keys, protected
health information, privileged legal material, or data they are not authorized
to share. The research team stores only the minimum artifact segments needed to
understand the incident. Screen-shared materials not approved for retention are
not captured in screenshots; the interviewer records a de-identified structural
summary. Names, company identifiers, exact monetary values, and unique workflow
details are generalized during transcription. The consent form separately
asks permission to quote de-identified speech and to reproduce redacted artifact
fragments. Every quotation used in the paper must be checked against the audio
or transcript and approved under the study's consent procedure.

\section{Complete Session Protocol}
\label{app:procedure}

\subsection{Pilot and freeze procedure}

Conduct two pilot sessions that are excluded from the reported sample. The
pilots test whether participants can understand the standardized scenario,
whether the open evidence-request phase yields actionable requests, whether the
common evidence fits within the allotted time, and whether the decision options
are interpreted distinctly. After the pilots, revise ambiguous wording and
then freeze: (1) the moderator script; (2) the standardized skill and incident;
(3) the evidence bank; (4) the common evidence set; (5) response-card wording;
and (6) the coding event schema. Any post-freeze change must be logged with its
reason, date, affected participants, and analytic consequence.

\subsection{Session overview}

\begin{table}[h]
\centering
\caption{Session structure and primary purpose.}
\label{tab:session-overview}
\begin{tabularx}{\linewidth}{P{0.13\linewidth}P{0.17\linewidth}Y P{0.18\linewidth}}
\toprule
\textbf{Time} & \textbf{Activity} & \textbf{Purpose} & \textbf{Primary data} \\
\midrule
0--5 min & Consent and orientation & Establish scope; emphasize that the task has no single expected diagnosis & Consent; baseline description \\
5--30 min & Real breakdown reconstruction & Examine current practice without imposing the triadic model & Timeline; artifact references; explanations; actions \\
30--65 min & Open-ended WoZ repair probe & Observe evidence seeking, scope revision, implementation choice, and action & Requests; predictions; proposals; final decision \\
65--80 min & Reflection and delegation interview & Compare probe and real practice; clarify desired authority & Interview; delegation sort; rationale \\
\bottomrule
\end{tabularx}
\end{table}

The nominal duration is 80 minutes. Small deviations are permitted only to let a participant finish a thought; the moderator must not selectively extend evidence seeking or decision time.
The facilitator should not reward a particular diagnosis, encourage a skill
edit, or imply that more evidence is always preferable.

\subsection{Opening script: 0--5 minutes}

The moderator reads:

\begin{quote}
We are studying how people understand and respond when a recurring AI workflow
does something they consider unacceptable. We are evaluating the workflow and
our study materials, not your skill. There may be more than one reasonable way
to interpret an incident, and you do not need to identify one correct cause.
During the second activity, you may request evidence from a prepared set. The
assistant will not have unlimited capabilities, and an unavailable request is
still useful data about what you wanted to know.
\end{quote}

The moderator confirms consent, recording preferences, artifact-retention
permissions, and the participant's preferred terminology for the reusable
configuration.

\subsection{Part 1: Artifact-grounded breakdown reconstruction, 5--30 minutes}

The moderator opens a neutral event timeline with six headings: \emph{Goal},
\emph{What happened}, \emph{What I observed}, \emph{What I inferred},
\emph{What I considered changing}, and \emph{What remained uncertain}. The
headings are presented as note-taking aids, not as a model of the system.

\paragraph{A. Establish the recurring workflow.}
\begin{enumerate}[leftmargin=*]
    \item What recurring work was the AI meant to perform?
    \item What makes two uses instances of the same recurring workflow?
    \item What outcomes, intermediate decisions, tool actions, or confirmation
    points matter to you?
    \item Who is accountable when the workflow behaves incorrectly?
\end{enumerate}

\paragraph{B. Reconstruct the incident chronologically.}
\begin{enumerate}[leftmargin=*]
    \item What did you ask the system to do, and what context was available?
    \item What did you directly observe in the output, trace, or tool results?
    \item What do you believe happened but did not directly observe?
    \item At what point did the workflow depart from your expectation?
    \item What consequence made the outcome unacceptable?
\end{enumerate}

The moderator asks the participant to mark each item as \emph{observed},
\emph{inferred}, or \emph{uncertain}. The moderator does not correct the
participant's account or introduce instruction-level terminology unless the
participant does so.

\paragraph{C. Elicit possible objects of repair.}
\begin{enumerate}[leftmargin=*]
    \item At the time, what did you think might need to change?
    \item What alternative changes or non-changes did you consider?
    \item What evidence supported each possibility?
    \item Were there explanations you could not distinguish?
    \item Was there any reason not to edit the reusable instructions?
\end{enumerate}
Only after the participant has answered these open prompts may the moderator use
a neutral rescue prompt: ``Did you consider changing anything about the desired
behavior, the persistent configuration, the current input, a tool, the amount
of automation, or something else?'' The transcript and coding record whether
an object of repair was raised spontaneously or after this prompt.

\paragraph{D. Reconstruct the actual response and release judgment.}
\begin{enumerate}[leftmargin=*]
    \item What did you or someone else actually change?
    \item Did you test only the motivating case or other cases as well?
    \item What behavior did you want to remain dependable?
    \item What made you willing to keep using, deploy, roll back, or abandon the
    revised workflow?
    \item Looking back, what information was missing?
\end{enumerate}

At the end of Part 1, the moderator writes a neutral one-sentence summary and
asks the participant to correct it. The summary must not identify a root cause
unless the participant has done so.

\subsection{Part 2: Open-ended Wizard-of-Oz repair probe, 30--65 minutes}

\subsubsection{Stage A: Incident only, 30--36 minutes}

Show only the materials in Appendix~\ref{app:scenario-incident}: the travel
request, observable tool calls, candidate flights, agent decision, and missed
presentation. Do not show the skill or an explicit list of requirements.

Ask:

\begin{enumerate}[leftmargin=*]
    \item What, if anything, is unacceptable about this execution?
    \item What should have happened in this case?
    \item What might you change first, if anything?
    \item What other future behavior could that change affect?
    \item How confident are you in your current explanation and action? Record
    each from 0 (not at all confident) to 100 (completely confident).
\end{enumerate}

Do not use the labels \emph{change}, \emph{preserve}, or \emph{unresolved} at
this stage.

\subsubsection{Stage B: Reveal the reusable artifact, 36--42 minutes}

Reveal the proposed six-instruction skill in Appendix~\ref{app:scenario-skill}
and the available tools. Give the participant time to read, then ask:

\begin{enumerate}[leftmargin=*]
    \item What does the skill change about your explanation, if anything?
    \item Which relationships does the skill state explicitly, and which are
    your inferences?
    \item Does the skill fully express the policy you think the workflow should
    follow?
    \item Would you change your proposed action?
    \item Update explanation and action confidence from 0--100.
\end{enumerate}

The participant then predicts the workflow's decision in the three short
contexts in Appendix~\ref{app:prediction-items} without seeing their outputs.
These predictions elicit the participant's current mental model; they are not
used to make inferential performance claims in the formative study.

\subsubsection{Stage C: Free evidence request, 42--52 minutes}

The moderator says:

\begin{quote}
You may now ask for any evidence that would help you decide what, if anything,
should change. You can ask to see another type of case, rerun a case, inspect a
tool result, try a what-if change, compare two revisions, or request something
else. Tell us what you hope each request will help you decide.
\end{quote}

For each request, the facilitator records:

\begin{itemize}[leftmargin=*]
    \item the participant's wording;
    \item the uncertainty or decision the request is intended to address;
    \item the matched evidence-bank card or an ``unavailable'' response;
    \item the participant's interpretation after seeing the evidence;
    \item any change in proposed object of repair, scope, confidence, or action.
\end{itemize}

The facilitator follows the lookup protocol in Appendix~\ref{app:wizard}. No
new model output, explanation, or diagnosis may be generated during the
session.

\subsubsection{Stage D: Common evidence, 52--59 minutes}

Present four evidence cards to every participant, using the counterbalanced
orders in Table~\ref{tab:common-orders}:

\begin{enumerate}[leftmargin=*]
    \item a routine preservation context;
    \item a difficult policy-boundary context;
    \item repeated executions of the motivating case;
    \item a context in which calendar/tool state is incomplete.
\end{enumerate}

If a participant already requested one of these cards, mark it as previously
seen and spend no additional time on it. After each new card, ask:

\begin{enumerate}[leftmargin=*]
    \item What, if anything, does this change about your explanation?
    \item Does it change what you would modify or where the modification should
    apply?
    \item Does it reveal behavior that should remain stable?
    \item Does it reveal a trade-off you are not ready to settle?
\end{enumerate}

Only after these open questions does the moderator introduce a worksheet with
three columns: \emph{Should behave differently}, \emph{Should remain stable},
and \emph{Not ready to decide}. The participant may leave any column empty.

\subsubsection{Stage E: Implementation and action choice, 59--65 minutes}

Ask the participant to select one implementation mode:

\begin{enumerate}[leftmargin=*]
    \item write a skill revision manually;
    \item ask AI to draft a revision and edit or approve it;
    \item describe the desired behavioral policy and let AI implement the text;
    \item defer implementation pending more evidence;
    \item retain the current Skill.
\end{enumerate}

If the participant requests an AI draft, provide a precomputed candidate that
corresponds to their selected policy class, not a live personalized generation.
The participant may annotate or reject it. Then show the available old-versus-
new evidence cards and ask for one final action:

\begin{itemize}[leftmargin=*]
    \item \textbf{Release}: reuse the revision under the reviewed conditions;
    \item \textbf{Revise}: change the candidate before reuse;
    \item \textbf{Gather more evidence}: withhold a decision; or
    \item \textbf{Defer}: retain the current Skill.
\end{itemize}

Record the participant's rationale and the evidence they consider decisive or
missing. The moderator must not imply that release is the desired endpoint.

\subsection{Part 3: Reflection and delegation interview, 65--80 minutes}

\paragraph{Compare real and standardized incidents.}
\begin{enumerate}[leftmargin=*]
    \item What felt similar or different between this scenario and your own
    workflow?
    \item Did the probe cause you to reconsider how you handled your real
    incident?
    \item Which evidence was useful, distracting, or missing?
\end{enumerate}

\paragraph{Representation and legibility.}
\begin{enumerate}[leftmargin=*]
    \item Would you rather begin from behavior, the skill text, or an execution
    trace? Why?
    \item What should be shown by default, and what should require drill-down?
    \item Do you need a responsible-instruction explanation, or is it enough to
    know how a candidate revision changes behavior?
    \item How should unstable or inconclusive evidence be represented?
\end{enumerate}

\paragraph{Division of labor.}
The participant sorts the responsibilities in
Table~\ref{tab:delegation-cards} into \emph{User}, \emph{AI}, or \emph{Shared},
then discusses any responsibility whose allocation depends on risk or artifact
familiarity.

\paragraph{Release and maintenance.}
\begin{enumerate}[leftmargin=*]
    \item What does ``remaining in control'' mean in this workflow?
    \item What evidence would be required before release in your own domain?
    \item What should be saved for the next maintainer or future revision?
    \item When should the system recommend no skill edit?
\end{enumerate}

\subsection{Debrief}

Explain that the standardized scenario was designed to admit multiple plausible
repairs and that the study did not contain a hidden correct diagnosis. Clarify
which outputs were precomputed, remind the participant not to adopt the study's
candidate patches in a real workflow without independent validation, and
confirm whether submitted materials may be retained.

\section{Standardized Travel-Recovery Scenario}
\label{app:scenario}

\DraftNote{The materials in this appendix are a proposed study instrument based on the paper's motivating travel-recovery scenario. They are not copied from a validated dataset. Pilot, audit, and freeze all wording and executions before data collection.}

\subsection{Motivating incident}
\label{app:scenario-incident}

A researcher's original flight is canceled at 14:00 on the day before a
conference presentation. The agent has access to the traveler's itinerary,
calendar, airline search, fare details, and rebooking tool. The observable
calendar result contains a paper presentation at 09:00 the following morning.
The search tool returns the two feasible options in
Table~\ref{tab:incident-options}.

\begin{table*}[t]
\centering
\caption{Candidate flights in the motivating incident.}
\label{tab:incident-options}
\begin{tabularx}{\textwidth}{P{0.08\textwidth}P{0.17\textwidth}P{0.16\textwidth}P{0.09\textwidth}P{0.11\textwidth}Y}
\toprule
\textbf{Option} & \textbf{Arrival} & \textbf{Airline} & \textbf{Fare} & \textbf{Connections} & \textbf{Consequence} \\
\midrule
A & 13:30 next day & Original airline & \$80 & 0 & Misses the 09:00 presentation \\
B & 20:00 same day & Different airline & \$420 & 1 & Arrives the evening before the presentation \\
\bottomrule
\end{tabularx}
\end{table*}

The agent selects Option A and explains that it is cheaper, uses the original
airline, and avoids a connection. It does not ask for confirmation. The study
shows the participant only observable inputs, tool outputs, and the final
decision; it does not reveal hidden chain-of-thought or claim that one internal
reason caused the decision.

\subsection{Proposed reusable skill}
\label{app:scenario-skill}

The skill is intentionally plausible but incomplete: it asks the agent to
retrieve relevant information and balance several preferences without fully
specifying how time-critical commitments override them.

\begin{enumerate}[leftmargin=*]
    \item Identify all disrupted flight segments and retrieve feasible
    alternatives for the traveler.
    \item Check the traveler's calendar for commitments within 24 hours of the
    original arrival time and include them in the comparison.
    \item Prefer the original airline when a feasible option is available.
    \item Prefer lower total fare and fewer connections among feasible options.
    \item Ask for confirmation before selecting an option priced above \$500 or
    before a non-refundable purchase.
    \item When no option satisfies all preferences, choose the option that
    minimizes overall travel disruption and explain the trade-off before
    rebooking.
\end{enumerate}

The six instructions create several plausible objects of repair without making
one a study-defined ground truth: the calendar information may be insufficiently
prioritized; airline or cost preferences may be overly broad; ``overall
travel disruption'' may be underspecified; confirmation may be too narrowly
triggered; or a tool/context failure may make a textual edit inappropriate.

\subsection{Prediction items after skill reveal}
\label{app:prediction-items}

Before any additional evidence, ask participants to predict the original
skill's likely choice and provide 0--100 confidence for each context:

\begin{enumerate}[leftmargin=*]
    \item \textbf{Routine leisure trip.} No calendar commitment depends on
    early arrival. A \$120 same-airline nonstop arrives at 15:00 next day; a
    \$840 different-airline one-stop arrives at 21:00 the same day.
    \item \textbf{Moderate commitment, aligned preferences.} A \$260
    same-airline nonstop arrives at 20:00 before a 10:00 meeting; a \$90
    same-airline flight arrives at 11:30 after the meeting begins.
    \item \textbf{High-cost boundary.} A \$1,180 different-airline flight
    arrives before a fixed 08:00 client presentation; a \$95 same-airline
    flight arrives six hours after it begins.
\end{enumerate}

The formative study uses these predictions to elicit participants' current
interpretation, not to estimate population-level prediction accuracy.

\subsection{Common evidence cards}
\label{app:common-evidence}

\paragraph{E1: Routine preservation context.}
Use the routine leisure trip above. The original skill's precomputed executions
should predominantly favor the lower-cost same-airline option. The evidence
card should expose raw decisions and explanations across repetitions, without
labeling this behavior as correct.

\paragraph{E2: Policy-boundary context.}
Use the high-cost boundary above. Precomputed outputs should illustrate the
trade-off between a fixed commitment and a large fare difference. It is
acceptable for outputs to vary, provided that the card accurately reports the
variation and does not manufacture a clean effect.

\paragraph{E3: Repeated motivating executions.}
Show five executions of the motivating incident under identical visible inputs.
Report each option selected, whether confirmation was requested, and the
observable justification. Do not average away disagreement.

\paragraph{E4: Incomplete calendar input.}
Show a separate execution in which the calendar tool returns an incomplete
result because the event is stored in a secondary calendar that is not queried.
This card tests how the Skill behaves when an instruction receives incomplete
tool-derived input.

\subsection{Candidate implementation classes}
\label{app:candidate-patches}

Candidate text is shown only when requested in Stage E. The study should include
at least these classes. Each is labeled as a candidate, with no recommended
status:

\paragraph{P1: Broad arrival-priority patch.}
``Always select the option with the earliest arrival, even when it costs more or
uses a different airline.'' This patch fixes the motivating case but is
intentionally likely to overreach in routine travel.

\paragraph{P2: Scoped commitment patch.}
``When the calendar contains a fixed commitment within 16 hours of arrival,
select an option scheduled to arrive at least three hours before that commitment.
Otherwise, continue to balance fare, airline, and connections. Ask for
confirmation whenever satisfying the commitment requires a fare more than
\$500 above the lowest feasible option.'' This patch encodes one possible
boundary and confirmation policy; the study does not treat its thresholds as
normatively correct.

\paragraph{P3: Uncertainty-and-confirmation patch.}
``If calendar information is missing, conflicting, or does not establish
whether arrival time is consequential, present the trade-off and ask the user
before rebooking.'' This patch leaves the priority open under uncertainty and
calls for confirmation.

\paragraph{P4: Defer the Skill edit.}
Retain the current Skill until the participant can state a supported behavioral
scope. This option records that the available evidence does not yet justify a
persistent revision.

\section{Wizard-of-Oz Evidence Bank and Facilitation Rules}
\label{app:wizard}

\subsection{Evidence-bank construction}

Freeze the evidence bank before the first analyzed session. Use
\ModelVersion{} accessed or frozen on \ModelDate{}, with the same system
configuration, tool outputs, sampling parameters, and execution harness for all
participants. Record model identifier, endpoint, temperature, maximum tokens,
tool schemas, prompt templates, and any unavailable seed control. Preserve raw
inputs, tool events, outputs, timestamps, and evaluator annotations.

At minimum, precompute:

\begin{itemize}[leftmargin=*]
    \item five repeated executions for the original skill in each scenario
    context;
    \item five executions for each candidate patch in the motivating,
    preservation, and boundary contexts;
    \item selected rule-level what-ifs, such as omitting the airline preference,
    fare preference, calendar check, or a relevant pair;
    \item complete and incomplete calendar-tool states;
    \item old-versus-new comparisons that expose both local success and possible
    behavioral overreach.
\end{itemize}

These runs support the formative probe; they do not substitute for the later
technical validation of any intervention method. If the precomputed model does
not naturally produce the intended variation, report the actual evidence.
Revise the study materials during the pilot; never manually edit model outputs.

\subsection{Evidence-card format}

Every evidence card contains:

\begin{enumerate}[leftmargin=*]
    \item the participant's request or the common question addressed;
    \item the exact context and configuration compared;
    \item a compact summary of decisions and observable actions;
    \item the number of repeated executions;
    \item a visible indicator of agreement or variation;
    \item expandable raw outputs and tool events;
    \item a neutral limitation statement, such as ``This comparison does not
    establish a unique cause.''
\end{enumerate}

The card must not contain a generated root-cause explanation, a responsible-
instruction label, or language that endorses a candidate patch.

\subsection{Request taxonomy and lookup}

The taxonomy in Table~\ref{tab:request-taxonomy} is used only by the facilitator
to retrieve a card; it is not shown to participants.

\begin{table}[h]
\centering
\caption{Wizard request taxonomy.}
\label{tab:request-taxonomy}
\begin{tabularx}{\linewidth}{P{0.22\linewidth}Y Y}
\toprule
\textbf{Request type} & \textbf{Examples} & \textbf{Response rule} \\
\midrule
Context request & ``Show a case without a hard commitment'' & Return nearest precomputed context; state differences explicitly \\
Repeat request & ``Run the same case again'' & Return all remaining repetitions, not a selected favorable subset \\
Artifact what-if & ``Ignore the airline rule'' & Return corresponding fixed intervention card \\
Tool-state request & ``Did the calendar return the event?'' & Return exact observable tool result; do not infer hidden state \\
Patch comparison & ``Compare broad versus scoped edits'' & Return fixed side-by-side old/new evidence \\
Raw evidence & ``Show me the underlying outputs'' & Expand unabridged executions and tool calls \\
Unsupported request & Request has no frozen card & State that the evidence is unavailable; ask what decision it would inform; log the gap \\
\bottomrule
\end{tabularx}
\end{table}

When multiple cards could match, use a prewritten priority order based on
semantic proximity, not on which card best supports the paper's hypothesis. A
second researcher should audit a random sample of session-to-card matches.

\subsection{Common-evidence counterbalancing}

Use the four orders in Table~\ref{tab:common-orders} in rotation. If the final
sample is not divisible by four, assign the remaining participants to the least
used orders.

\begin{table}[h]
\centering
\caption{Counterbalanced order of common evidence cards.}
\label{tab:common-orders}
\begin{tabular}{ccccc}
\toprule
\textbf{Order} & \textbf{1st} & \textbf{2nd} & \textbf{3rd} & \textbf{4th} \\
\midrule
A & E1 & E2 & E3 & E4 \\
B & E2 & E3 & E4 & E1 \\
C & E3 & E4 & E1 & E2 \\
D & E4 & E1 & E2 & E3 \\
\bottomrule
\end{tabular}
\end{table}

\subsection{Facilitator neutrality rules}

The facilitator must:

\begin{itemize}[leftmargin=*]
    \item use the participant's terminology and delay introducing ``skill,''
    ``policy,'' or ``workflow'';
    \item ask what a request is intended to decide before returning evidence;
    \item return all relevant repeated runs, including discordant examples;
    \item avoid praising additional testing or portraying release as success;
    \item avoid stating that an intervention demonstrates causality;
    \item record unavailable requests and do not approximate them with a live
    model;
    \item disclose when a candidate patch was precomputed and not personalized;
    \item keep a deviation log for any departure from the script.
\end{itemize}

\section{Study Instruments}
\label{app:instruments}

\subsection{Breakdown timeline worksheet}

\begin{table}[h]
\centering
\caption{Neutral breakdown-reconstruction worksheet.}
\label{tab:timeline}
\small
\begin{tabularx}{\linewidth}{P{0.20\linewidth}Y P{0.22\linewidth}}
\toprule
\textbf{Prompt} & \textbf{Participant notes} & \textbf{Evidence status} \\
\midrule
Goal / expected outcome & \DataNeeded{SESSION NOTES} & Observed / inferred / uncertain \\
What happened & \DataNeeded{SESSION NOTES} & Observed / inferred / uncertain \\
Possible explanations & \DataNeeded{SESSION NOTES} & Evidence for / against / missing \\
Possible Skill explanation & \DataNeeded{SESSION NOTES} & Policy / instruction / interaction / activation / uncertain \\
Actions considered or taken & \DataNeeded{SESSION NOTES} & Implemented / rejected / deferred \\
What remained uncertain & \DataNeeded{SESSION NOTES} & Open question \\
\bottomrule
\end{tabularx}
\end{table}

The ``possible Skill explanation'' categories are completed by the researcher
after the interview, not shown to the participant during reconstruction.

\subsection{Stage-transition measures}

At the end of Stages A, B, C, and D, record:

\begin{enumerate}[leftmargin=*]
    \item current explanation or set of explanations;
    \item proposed action or non-action;
    \item confidence in the explanation (0--100);
    \item confidence in the proposed action's scope (0--100);
    \item behavior believed likely to change elsewhere;
    \item behavior believed important to retain;
    \item unresolved trade-offs;
    \item evidence still desired.
\end{enumerate}

These ratings support process tracing and within-participant comparison. Do not
use them to claim causal effects of individual evidence cards or population-
level calibration.

\subsection{Repair-scope worksheet}

After the common evidence set, use Table~\ref{tab:scope-worksheet}. Ask the
participant to phrase entries behaviorally, not as instruction numbers.

\begin{table}[h]
\centering
\caption{Repair-scope worksheet introduced only after open evidence seeking.}
\label{tab:scope-worksheet}
\begin{tabularx}{\linewidth}{YYY}
\toprule
\textbf{Should behave differently} & \textbf{Should remain stable} & \textbf{Not ready to decide} \\
\midrule
\DataNeeded{PARTICIPANT ENTRY} & \DataNeeded{PARTICIPANT ENTRY} & \DataNeeded{PARTICIPANT ENTRY} \\
\bottomrule
\end{tabularx}
\end{table}

For each entry, ask which concrete context supports it and whether the judgment
reflects a domain policy, an assumption about the current skill, or an observed
execution.

\subsection{Delegation cards}

\begin{table}[h]
\centering
\caption{Responsibilities sorted into User, AI, or Shared during the exit interview.}
\label{tab:delegation-cards}
\begin{tabularx}{\linewidth}{P{0.34\linewidth}Y}
\toprule
\textbf{Responsibility} & \textbf{Follow-up prompt} \\
\midrule
Decide whether the observed behavior is acceptable & What knowledge or accountability is required? \\
Generate contrasting future contexts & When should AI ask before generating cases? \\
Repeat executions and summarize variation & How much raw evidence is needed? \\
Choose a skill-level what-if intervention & Should the user approve each intervention? \\
Choose an instruction-level what-if or interaction test & Which instructions should remain fixed? \\
Draft a textual patch & Must the user understand every textual change? \\
Edit or approve the final skill text & Does this depend on artifact familiarity or stakes? \\
Judge whether the candidate changed the intended behavior & What context coverage is enough? \\
Decide whether established behavior was preserved & Who defines preservation commitments? \\
Release, revise, gather evidence, or defer & Who holds final authority and why? \\
Save rationale and regression cases for future maintenance & What should become persistent documentation? \\
\bottomrule
\end{tabularx}
\end{table}

\subsection{Final action record}

For the selected action, record:

\begin{itemize}[leftmargin=*]
    \item implementation mode;
    \item final action category;
    \item evidence considered sufficient;
    \item evidence considered missing;
    \item accepted unresolved regions;
    \item perceived consequences of a wrong decision;
    \item who should be notified or involved in a real deployment;
    \item what should be preserved as a future regression case.
\end{itemize}

\section{Data Analysis and Claim Audit}
\label{app:analysis}

\subsection{Analytic units and provenance}

Segment the corpus into \emph{events}: a participant states an expectation,
reports an observation, offers an explanation, requests or receives evidence,
changes an interpretation, revises repair scope, delegates a responsibility, or
makes an action decision. Every event retains links to participant, activity
(real incident or probe), timestamp, artifact/evidence card, and preceding
state. This provenance prevents quotations from being separated from the
evidence conditions under which they occurred.

A main-text finding should not rest solely on stated preferences. Unless the
finding is explicitly about current practice, support it with at least two data
forms, such as a real-incident account plus a probe action, or a participant
request plus a later interview rationale.

\subsection{First-cycle inductive event coding}

The first cycle does not begin with the user--skill--workflow triad. Analysts
code participants' own terms using the families in Table~\ref{tab:first-cycle}.
The listed codes are starting prompts, not a closed codebook.

\begin{table*}[t]
\centering
\caption{Initial event-code families.}
\label{tab:first-cycle}
\begin{tabularx}{\textwidth}{P{0.18\textwidth}P{0.29\textwidth}Y}
\toprule
\textbf{Family} & \textbf{Example codes} & \textbf{Analytic question} \\
\midrule
Expected behavior & goal, constraint, priority, exception, confirmation & What did the participant believe should happen? \\
Observed breakdown & output symptom, tool action, omitted step, consequence & What was directly observable? \\
Inference & presumed priority, hidden interpretation, guessed tool state & Which elements came from observation, and which from inference? \\
Candidate explanation & policy gap, instruction, interaction, activation condition, variability, tool-derived input & How might the Skill have produced the behavior? \\
Evidence & artifact text, raw run, repetition, contrastive case, what-if, version history & What supported or weakened an explanation? \\
Uncertainty & competing explanation, unstable behavior, unresolved policy, missing evidence & What remained unsettled? \\
Repair action & textual edit, AI-drafted edit, rollback, defer, no action & What Skill change was considered or taken? \\
Scope move & narrow, broaden, reframe, preserve, mark unresolved, abandon & How did the intended reach change? \\
Delegation & user, AI, shared; evidence, implementation, policy, release & Who performed or owned the work? \\
Decision & release, revise, gather evidence, defer & What action followed, and why? \\
\bottomrule
\end{tabularx}
\end{table*}

Analysts write a memo after each session identifying emergent codes, surprising
requests, and evidence that contradicts the current story. The research team
reviews these memos before modifying the codebook.

\subsection{FQ1: Abductive instruction--behavior synthesis}

After first-cycle coding, examine how participants connected observed behavior
to the reusable Skill. Candidate second-cycle categories may include:

\begin{itemize}[leftmargin=*]
    \item intended policy or domain judgment;
    \item individual Skill instruction;
    \item interaction among Skill instructions;
    \item context-dependent instruction activation;
    \item variable enacted behavior;
    \item insufficient evidence for a Skill edit.
\end{itemize}

The triadic account is retained only if intended behavior, reusable Skill, and
enacted workflow form a recurring and explanatory structure across cases.
Compare participants with different artifact familiarity to determine whether
instruction--behavior uncertainty reflects implementation expertise or a
broader property of Skill-mediated workflow control.

\subsection{FQ2: Evidence-to-action process tracing}

Construct a trajectory table for every participant:

\begin{table*}[t]
\centering
\caption{Within-participant process-trace schema.}
\label{tab:process-trace}
\begin{tabularx}{\textwidth}{P{0.06\textwidth}P{0.17\textwidth}P{0.20\textwidth}P{0.20\textwidth}P{0.14\textwidth}Y}
\toprule
\textbf{Time} & \textbf{Evidence requested / received} & \textbf{Interpretation before $\rightarrow$ after} & \textbf{Object of repair before $\rightarrow$ after} & \textbf{Scope move} & \textbf{Action / confidence move} \\
\midrule
\DataNeeded{TIME} & \DataNeeded{EVENT} & \DataNeeded{CHANGE} & \DataNeeded{CHANGE} & \DataNeeded{MOVE} & \DataNeeded{MOVE} \\
\bottomrule
\end{tabularx}
\end{table*}

Distinguish three effects:

\begin{enumerate}[leftmargin=*]
    \item \emph{confirmatory}: evidence changes confidence but not
    interpretation or scope;
    \item \emph{specification-changing}: evidence adds, removes, or changes a
    behavioral commitment or policy boundary;
    \item \emph{mapping-changing}: evidence changes which instruction or
    interaction the participant considers relevant to the repair.
\end{enumerate}

Compare self-requested and common evidence. A claim that contrastive cases help
construct repair scope requires evidence of specification-changing transitions.
A confidence increase after a researcher-provided example is insufficient.

\subsection{FQ3: Enacted and stated delegation}

Create a matrix with rows for evidence work, implementation work, policy
judgment, and release decision, and columns for:

\begin{itemize}[leftmargin=*]
    \item what the participant requested AI to do in the probe;
    \item what implementation mode they selected;
    \item how they allocated the corresponding delegation card;
    \item what they said they would prefer in their own workflow;
    \item conditions that changed the allocation, such as consequence,
    familiarity, traceability, or uncertainty.
\end{itemize}

Code disagreement among these forms as data. Do not privilege interview
statements. For example, a participant may claim to prefer manual
control but choose AI drafting when facing time pressure, or may delegate text
editing while retaining a release veto.

\subsection{Researcher process and reflexivity}

Two researchers independently analyze an initial subset selected to maximize
variation in artifact familiarity, workflow consequence, and ownership
arrangement. They compare event boundaries, identify missing code families, and
produce a shared event schema. One researcher then leads full-corpus coding;
the second reviews at least 30\% of sessions, all candidate negative cases, and
all excerpts considered for the paper. Disagreement is used to refine theme
boundaries. A reliability statistic would be inconsistent with this
interpretive analysis.

Maintain:

\begin{itemize}[leftmargin=*]
    \item a decision log linking codebook and theme revisions to evidence;
    \item a deviation log for study-protocol changes;
    \item a quotation audit with source timestamp and consent status;
    \item a finding matrix showing supporting and disconfirming cases;
    \item separate memos for real incidents and standardized probe behavior
    before cross-comparison;
    \item a reflexivity memo describing the team's assumptions about skills,
    repair, user control, and the proposed system.
\end{itemize}

\subsection{Falsification and claim thresholds}

Before finalizing findings, explicitly test the propositions in
Table~\ref{tab:falsification}. This check protects against a study that can only
confirm the intended system story; it does not constitute statistical
falsification.

\begin{table}[h]
\centering
\caption{Evidence that would strengthen, weaken, or redirect the paper's core claims.}
\label{tab:falsification}
\begin{tabularx}{\linewidth}{P{0.24\linewidth}YY}
\toprule
\textbf{Proposition} & \textbf{Evidence that supports it} & \textbf{Evidence that weakens or redirects it} \\
\midrule
Skill text does not fully expose instruction effects & Participants revise which instructions or interactions they consider relevant after execution evidence; experts also report context-dependent uncertainty & Most participants reliably identify the relevant instruction and predict behavior from text alone \\
Repair scope is constructed across cases & Evidence produces new preservation commitments, changed policy boundaries, unresolved regions, or no-edit decisions & Participants begin with complete policies; additional cases only confirm them \\
Control centers on policy and release, not necessarily manual text authorship & Participants delegate evidence or patch drafting while retaining behavior and release decisions & Participants either insist on complete manual textual ownership or willingly delegate policy and release without review \\
Behavioral patch review is needed beyond ordinary test-and-diff & Participants revise the object or specification of repair after comparing fixed versions & The main difficulty is simply insufficient regression coverage after a known patch \\
\bottomrule
\end{tabularx}
\end{table}

A main finding should satisfy four thresholds:

\begin{enumerate}[leftmargin=*]
    \item it is supported by multiple participants and more than one data form;
    \item its boundary or heterogeneity is reported;
    \item at least one plausible disconfirming case has been sought and
    addressed;
    \item the wording does not imply causal effectiveness that only a later
    controlled evaluation can establish.
\end{enumerate}

If the evidence does not support the provisional finding titles in
Section~\ref{sec:formative-findings}, replace them. In particular, use the
controlled study to claim improved workflow legibility, calibrated prediction,
reduced repair-scope mismatch, or better release decisions. Formative interviews
alone cannot establish these effects.

\subsection{Descriptive reporting}

Permissible descriptive summaries include:

\begin{itemize}[leftmargin=*]
    \item number of incidents and candidate objects of repair;
    \item frequencies of spontaneously requested evidence types;
    \item numbers of narrow, broaden, reframe, defer, abandon, and no-edit
    transitions;
    \item implementation and final-action choices;
    \item counts of participants contributing to a theme.
\end{itemize}

Report denominators, distinguish real-incident and probe counts, and avoid
phrases such as ``most users'' without the sample denominator. Confidence
ratings and behavioral predictions are elicitation aids; unless the study is
redesigned and powered accordingly, do not test them inferentially.

\section{Material Freeze and Reproducibility Checklist}
\label{app:freeze}

Before the first analyzed session, archive the following with immutable version
identifiers:

\begin{enumerate}[leftmargin=*]
    \item recruitment screener, consent form, and artifact request;
    \item moderator script and participant worksheets;
    \item standardized incident, skill, contexts, and candidate patch set;
    \item all raw evidence-bank inputs and outputs;
    \item model, parameter, tool, and harness configuration;
    \item evidence cards and request-to-card lookup table;
    \item common-evidence counterbalancing schedule;
    \item pilot changes and the final freeze decision;
    \item initial event schema, analytic memos, and deviation-log templates.
\end{enumerate}

After data collection, archive:

\begin{enumerate}[leftmargin=*]
    \item de-identified transcripts and retained artifact fragments, subject to
    consent;
    \item logs linking sessions to evidence cards and counterbalancing assignments;
    \item coded event data and process traces;
    \item codebook and theme revision history;
    \item finding-by-participant and negative-case matrices;
    \item quotation audit and consent status;
    \item the exact counts and excerpts used in the manuscript.
\end{enumerate}

If public release of participant data is not possible, release the standardized
materials, evidence bank, moderator script, codebook, synthetic examples of the
analysis tables, and a data-access statement explaining restrictions.

\DraftNote{Before submission, remove all draft notes and red placeholders; replace the provisional process figure caption; report the actual sample, ethics status, model version, pilot changes, deviations, and negative cases; and ensure that every finding is traceable to audited data.}
